% Part: first-order-logic
% Chapter: natural-deduction
% Section: provability-quantifiers

% verification of properties of provability needed for maximally
% consistent sets in the completeness chapter.

\documentclass[../../../include/open-logic-section]{subfiles}

\begin{document}

\olfileid{fol}{ntd}{qpr}

\olsection{\usetoken{S}{derivability} والمكممات}

\begin{explain}
  ومن مقدمات برهان الاكتمال أن تثبت بقواعد الاستنتاج الطبيعي
  النتائج الآتية في علاقة~$\Proves$؛ فلذلك نفردها بالبيان هنا.
\end{explain}

\begin{thm}
\ollabel{thm:strong-generalization} إذا كان~$c$ رمزًا ثابتًا لا يرد
في~$\Gamma$ ولا في~$!A(x)$، وكان~$\Gamma \Proves !A(c)$، فإن
$\Gamma \Proves \lforall[x][!A(x)]$.
\end{thm}

\begin{proof}
نأخذ~$\delta$، وهو !!a{derivation} لـ~$!A(c)$ من~$\Gamma$، ونختمه
باستدلال \Intro{\lforall}؛ فيتم بذلك !!a{derivation} لـ
$\lforall[x][!A(x)]$. وهذا جائز، إذ~$c$ لا يرد في~$\Gamma$ ولا
في~$!A(x)$؛ فلا يختل شرط المتغير المميَّز.
\end{proof}

\begin{prop}
\ollabel{prop:provability-quantifiers}
\begin{tagenumerate}{prvEx,prvAll}
\tagitem{prvEx}{$!A(t) \Proves \lexists[x][!A(x)]$.}{}

\tagitem{prvAll}{$\lforall[x][!A(x)] \Proves !A(t)$.}{}
\end{tagenumerate}
\end{prop}

\begin{proof}
\begin{tagenumerate}{prvEx,prvAll}
\tagitem{prvEx}{يكفي !!a{derivation} لـ~$\lexists[x][!A(x)]$ من~$!A(t)$
  بإدخال المكمّم الوجودي مرة واحدة:
\begin{prooftree}
\AxiomC{$!A(t)$}
\RightLabel{\Intro{\lexists}}
\UnaryInfC{$\lexists[x][!A(x)]$}
\end{prooftree}}{}

\tagitem{prvAll}{وأما !!a{derivation} لـ~$!A(t)$ من
  $\lforall[x][!A(x)]$ فيكفي فيه حذف المكمّم الكلي مرة واحدة:
\begin{prooftree}
\AxiomC{$\lforall[x][!A(x)]$}
\RightLabel{\Elim{\lforall}}
\UnaryInfC{$!A(t)$}
\end{prooftree}}{}
\end{tagenumerate}
\end{proof}

\end{document}
